\topskip0pt
\vspace*{\fill}
\begin{adjustwidth}{0.5cm}{0.5cm}
	
	\begin{center}
		\textbf{\Large The Hitchhiker's Guide to Applied Cryptography.}
	\end{center}
	
	~\\
	\textbf{What this handbook is for.} This handbook provides a \textit{guided overview of relevant and current concepts in applied cryptography}. Each topic is accompanied by a brief explanation and, where appropriate, pointers to additional resources for deeper study. In other words, the handbook is a \textit{starting point}, curated based on what we consider essential knowledge in this context, without claiming to be exhaustive or fully comprehensive. \underline{Please check out the latest} \underline{version of this handbook at \href{https://github.com/aleph-fbk/hgac}{\textbf{github.com/aleph-fbk/hgac}}}. \\
	
	\textbf{What this handbook is not.} The handbook is \textit{neither} a textbook, nor an encyclopedia, nor a dictionary. Similarly, the handbook is by no means a horizontally complete or vertically exhaustive treatment of all facets of (especially theoretical) cryptography. In particular, mathematical details, formal proofs, and descriptions of the inner functioning of cryptographic algorithms are intentionally minimized (where not omitted entirely) to maintain brevity and focus. Instead, the emphasis is on practical understanding and conceptual clarity, while pointers to authoritative sources are available for readers who wish to explore further (see \Cref{sec:appendix:further_resources}). \\
	
	\textbf{How to read this handbook.} The handbook is \textit{structured as a series of self-contained notes}. Each note can be read independently, though references to related notes (denoted with ``\texttt{§}'') or more detailed sources are included where useful. While a top-down reading of the handbook is possible, readers can \textit{read any one single note} based on their interests or needs. While prior familiarity with cryptography is helpful, this handbook strives to remain accessible to readers with a general background in computer science or cyber-security. \\
	
	%\textbf{Intended audience.} This handbook is primarily aimed at anyone who ($i$) requires a practical understanding of cryptography without delving (too much) into its theoretical underpinnings, ($ii$) seeks a reference companion to standard cryptographic texts, or ($iii$) needs a high-level overview of various cryptographic components and their relationships. \\
	
	\textbf{How to contribute.} This is a living handbook curated by the Applied Cryptography (\textit{ALEPH}) research unit of the \textit{Center for Cybersecurity} in \textit{Fondazione Bruno Kessler}. Contributions such as corrections, suggestions, or improved formulations are welcome. Please refer to \href{https://github.com/aleph-fbk/hgac/blob/main/CONTRIBUTING.md}{\underline{\textbf{CONTRIBUTING.md}}} for more information. Finally, note that any references to existing cryptographic services, libraries, and hardware should not be interpreted as an endorsement of such services, libraries, and hardware (unless stated otherwise). \\

	\begin{center}
		{\small\textit{If you come across any content in this handbook that lacks proper attribution, we sincerely apologize. We have compiled material over many years, hence some sources may not have been documented accurately. Please let us know at} \ul{\texttt{hgac@fbk.eu}}\textit{, and we will promptly make the necessary corrections.}}
	\end{center}

\end{adjustwidth}
\vspace*{\fill}

\clearpage